\begin{rem}
When $\theta = 1$, Theorem~\ref{macdonaldthm1} recovers \cite[Theorem~3.6]{GP}.
\end{rem}

\begin{rem}\label{macdonaldthetaNinteger}
The inf\/inite contours in Theorems~\ref{macdonaldthm2} and~\ref{macdonaldthm25} are needed because there are inf\/initely many real poles in the integrand and the contour needs to enclose all of them. When $\theta\in\N$, the integrands have f\/initely many poles, and we can therefore close the contours, obtaining eventually Theorem~\ref{macdonaldthm1}. More generally, if $\theta > 0$ is such that $\theta N \in\N$, a similar remark applies. In fact, we can write the product of $q$-Gamma function ratios appearing in the integrand~\eqref{macdonaldthm2eqn} as
\begin{gather}
\prod_{i=1}^N{\frac{\Gamma_q(\lambda_i + \theta(N-i)-z)}{\Gamma_q(\lambda_i + \theta(N-i+1)-z)}}\nonumber\\
 \qquad{} = \frac{\Gamma_q(\lambda_1 + \theta(N-1) - z)}{\Gamma_q(\lambda_2 + \theta(N-1)-z)}\cdots\frac{\Gamma_q(\lambda_{N-1} + \theta - z)}{\Gamma_q(\lambda_N + \theta-z)}\frac{\Gamma_q(\lambda_N - z)}{\Gamma_q(\lambda_1 + \theta N - z)},\label{ratioqgammas}
\end{gather}
and since $\Gamma_q(t+1) = \frac{1 - q^t}{1 - q}\Gamma_q(t)$, we conclude that the product above is a rational function in $q^{-z}$ with f\/initely many real poles. Thus formula~\eqref{macdonaldthm2eqn} is true if we replaced contour~$\CC^+$ by a closed contour $\CC_0$ containing all f\/initely many real poles of the integrand. Similarly, we can replace $\CC^-$ by a closed contour $\CC_0$ in~\eqref{macdonaldthm25eqn}.
\end{rem}
